\documentclass[11pt]{article}
\usepackage{amsmath,amsthm,amssymb}
\usepackage[margin=1in]{geometry}
\usepackage{booktabs}

\newtheorem{theorem}{Theorem}[section]
\newtheorem{lemma}[theorem]{Lemma}
\newtheorem{proposition}[theorem]{Proposition}
\newtheorem{corollary}[theorem]{Corollary}
\theoremstyle{definition}
\newtheorem{definition}[theorem]{Definition}
\theoremstyle{remark}
\newtheorem{remark}[theorem]{Remark}

\newcommand{\ZZ}{\mathbb{Z}}
\newcommand{\CC}{\mathbb{C}}
\newcommand{\Sym}{\mathfrak{S}}
\newcommand{\im}{\operatorname{im}}
\newcommand{\rk}{\operatorname{rank}}
\DeclareMathOperator{\Ind}{Ind}
\DeclareMathOperator{\Res}{Res}

\title{Even-Block Gap Closure at $k = 4$\\for the Rank Isolation Lemma}
\author{Clio}
\date{April 17, 2026}

\begin{document}
\maketitle

\begin{abstract}
We prove that the even-block gap in the Rank Isolation Lemma closes
unconditionally at $k = 4$.  The Rank Isolation Lemma asserts that rank
drops in the staircase symmetrizing product
$\Pi = \prod (1 + s_{i_j})/2$ for the longest element $w_0 \in \Sym_n$
occur \emph{only} at first appearances of the $H$-generators
$s_1, s_3, s_5, \ldots$, where $H = \langle s_1, s_3, s_5, \ldots \rangle$.
The within-block cascade and odd block starts are proved elsewhere.
The even-block gap---whether $P_k$ preserves rank at the start of an
even block---reduces by cascade transport to a single step within each
block, and by parabolic reduction to $\Sym_{k+1}$ irreps.

We show that in any $\Sym_n$-representation ($n \ge 5$), the
intersection $E_1^{(+1)} \cap E_3^{(+1)} \cap
\ker\bigl((1{+}s_4)(1{+}s_2)/4\bigr)$ is trivial.  The proof uses two
applications of an \emph{adjacent disjointness} lemma for eigenspaces
of consecutive Coxeter generators.  We explain why the analogous argument
fails for $k \ge 6$.  The result has been verified computationally for
all irreps of $\Sym_n$, $n = 5, 6, 7$.
\end{abstract}


%% ====================================================================
\section{Context and setup}
%% ====================================================================

\subsection{The staircase product}

Let $\Sym_n$ act on a finite-dimensional complex vector space $V$
via simple transpositions $s_1, \ldots, s_{n-1}$ satisfying the usual
Coxeter relations.  Define the \emph{staircase word} for the longest
element $w_0$ to be
\[
  s_1,\; s_2 s_1,\; s_3 s_2 s_1,\; \ldots,\; s_{n-1} s_{n-2} \cdots s_1,
\]
and the \emph{staircase product}
\[
  \Pi \;=\; \prod_{j} \frac{1 + s_{i_j}}{2},
\]
where the factors are taken in the order prescribed by the staircase
word (left to right).  Each factor $P_i = (1 + s_i)/2$ is an
idempotent projection onto the $+1$ eigenspace $E_i^{(+1)}$ of~$s_i$.

\subsection{The $H$-invariant subspace}

Let $H = \langle s_1, s_3, s_5, \ldots \rangle \le \Sym_n$ be the
subgroup generated by odd-indexed simple transpositions.  The
\emph{Rank Isolation Lemma} (conjectured and verified computationally
through $n = 16$) states:

\begin{conjecture}[Rank Isolation]
For any irreducible $\Sym_n$-representation $V_\lambda$,
\[
  \rk\bigl(\Pi\big|_{V_\lambda}\bigr) \;=\; \dim V_\lambda^H.
\]
Moreover, the rank drops occur \emph{only} at the first appearance of
each $H$-generator $s_{2m-1}$ in the staircase word.
\end{conjecture}

Three of the four injectivity components are proved unconditionally:
\begin{enumerate}
  \item \textbf{Within-block cascade:} Adjacent disjointness gives
    $E_i^{(-1)} \cap E_{i+1}^{(+1)} = \{0\}$, making each step within
    a block injective.
  \item \textbf{Block~2 start:} $P_2$ on $E_1^{(+1)}$ is injective
    by adjacent disjointness.
  \item \textbf{Odd block starts:} For $s_{k} \in H$ (odd~$k$), the
    transported image already lies in $E_k^{(+1)}$, and the
    contraction argument proves injectivity.
\end{enumerate}

The remaining gap is at \emph{even} block starts: when block~$k$
begins with $P_k$ and $k$ is even, $k \ge 4$.

\subsection{Reduction of the even-block gap}

Two reductions narrow the gap:

\begin{itemize}
  \item \textbf{Cascade transport} (proved separately): The gap at the
    start of even block~$k$ reduces to a \emph{single step} within
    block~$k{-}1$, namely $P_{k-2}$ applied immediately after $P_{k-1}$.
  \item \textbf{Parabolic reduction} (proved separately): Even
    block~$k$ holds for all~$n$ if and only if it holds for all
    $\Sym_{k+1}$ irreps.
\end{itemize}

\noindent
Thus, closing even block $k = 4$ requires proving a statement about
$\Sym_5$ representations---which we do below for arbitrary $\Sym_n$.


%% ====================================================================
\section{Adjacent disjointness}
%% ====================================================================

The following elementary lemma is the engine of the proof.

\begin{lemma}[Adjacent disjointness]\label{lem:adj-disj}
Let $V$ be any representation of\/ $\Sym_n$ ($n \ge 3$), and let
$s_i, s_{i+1}$ be consecutive simple transpositions.  Then
\[
  E_i^{(+1)} \cap E_{i+1}^{(-1)} \;=\; \{0\}
  \qquad\text{and}\qquad
  E_i^{(-1)} \cap E_{i+1}^{(+1)} \;=\; \{0\},
\]
where $E_j^{(\pm 1)} = \{v \in V : s_j\, v = \pm v\}$.
\end{lemma}

\begin{proof}
The subgroup $\langle s_i, s_{i+1} \rangle \cong \Sym_3$ acts on~$V$
by restriction.  It suffices to verify the claim on each irreducible
$\Sym_3$-representation.

\begin{itemize}
  \item \textbf{Trivial representation} ($\dim = 1$):
    $s_i$ and $s_{i+1}$ both act as $+1$.
    So $E_i^{(+1)} = V$, $E_{i+1}^{(-1)} = \{0\}$,
    and $E_i^{(-1)} = \{0\}$, $E_{i+1}^{(+1)} = V$.
    Both intersections are trivial.

  \item \textbf{Sign representation} ($\dim = 1$):
    $s_i$ and $s_{i+1}$ both act as $-1$.
    So $E_i^{(+1)} = \{0\}$ and $E_{i+1}^{(+1)} = \{0\}$.
    Both intersections are trivial.

  \item \textbf{Standard representation} ($\dim = 2$):
    In a suitable basis, $s_i$ and $s_{i+1}$ are represented by
    reflection matrices with eigenvalues $\{+1, -1\}$.
    Each eigenspace is one-dimensional.  The two reflections
    correspond to distinct hyperplanes in $\mathbb{R}^2$
    (the angle between them is $\pi/3$, since
    $(s_i s_{i+1})^3 = 1$).  In particular,
    $E_i^{(+1)} \ne E_{i+1}^{(-1)}$ and
    $E_i^{(-1)} \ne E_{i+1}^{(+1)}$: a $+1$-eigenline of one
    reflection is not the $-1$-eigenline of the other.
    Since both eigenspaces are one-dimensional, distinct lines
    in $\mathbb{R}^2$ meet only at the origin.
\end{itemize}

\noindent
Since every $\Sym_3$-representation decomposes into these three
irreducibles, the result follows by taking direct sums.
\end{proof}


%% ====================================================================
\section{Even-block closure at $k = 4$}
%% ====================================================================

\begin{lemma}[Even-block closure at $k = 4$]\label{lem:k4}
Let $V$ be any representation of\/ $\Sym_n$ ($n \ge 5$).  For any
$v \in V$ satisfying $s_1\, v = v$ and $s_3\, v = v$, if
\[
  (1 + s_4)(1 + s_2)\, v \;=\; 0,
\]
then $v = 0$.

\noindent
Equivalently:
\[
  E_1^{(+1)} \;\cap\; E_3^{(+1)} \;\cap\;
  \ker\!\left(\frac{(1 + s_4)(1 + s_2)}{4}\right) \;=\; \{0\}.
\]
\end{lemma}

\begin{proof}
Suppose $(1 + s_4)(1 + s_2)\, v = 0$ with $s_1\, v = v$ and
$s_3\, v = v$.  Since $s_4$ is an involution, decompose
\[
  v \;=\; w_+ + w_-, \qquad s_4\, w_+ = w_+,\quad s_4\, w_- = -w_-.
\]

\medskip
\noindent\textbf{Step 1 (Kill $w_+$).}\;
Since $|1 - 4| = 3 \ge 2$, the transposition $s_1$ commutes with
$s_4$.  Therefore the $E_4^{(\pm 1)}$ decomposition is preserved
by~$s_1$, and in particular $s_1\, w_+ = w_+$ (projecting
$s_1\, v = v$ to $E_4^{(+1)}$).

Similarly, $|2 - 4| = 2 \ge 2$, so $s_2$ commutes with $s_4$.
Thus $(1 + s_2)$ preserves the $E_4^{(\pm 1)}$ decomposition.
Using the commutativity $[s_2, s_4] = 0$, we rewrite the hypothesis:
\[
  0 \;=\; (1 + s_4)(1 + s_2)\, v
    \;=\; (1 + s_2)(1 + s_4)\, v
    \;=\; (1 + s_2)(2 w_+).
\]
Hence $(1 + s_2)\, w_+ = 0$, which gives $s_2\, w_+ = -w_+$,
i.e., $w_+ \in E_2^{(-1)}$.

Combining: $w_+ \in E_1^{(+1)} \cap E_2^{(-1)}$.  By adjacent
disjointness (Lemma~\ref{lem:adj-disj} applied to $s_1, s_2$),
this intersection is trivial.  Therefore $\boldsymbol{w_+ = 0}$.

\medskip
\noindent\textbf{Step 2 (Kill $w_-$).}\;
Since $w_+ = 0$, we have $v = w_- \in E_4^{(-1)}$.  But by
hypothesis $s_3\, v = v$, so $v \in E_3^{(+1)}$.  Hence
$v \in E_3^{(+1)} \cap E_4^{(-1)}$.  By adjacent disjointness
(Lemma~\ref{lem:adj-disj} applied to $s_3, s_4$), this intersection
is trivial.  Therefore $\boldsymbol{v = 0}$.
\end{proof}

\begin{corollary}[Even block $k = 4$ is safe]\label{cor:k4-safe}
At the gap step for even block~$4$ in the staircase product, $P_2$
preserves the rank of the transported $V^H$ image.  Consequently, no
rank drop occurs at even block~$4$.
\end{corollary}

\begin{proof}
By cascade transport, the gap at block~$4$ reduces to a single step:
$P_2$ applied to vectors in the transported $V^H$ image after $P_3$.
The transported image at this point satisfies two eigenspace conditions:

\begin{enumerate}
  \item $s_3\, w = w$: since $P_3 = (1 + s_3)/2$ projects onto
    $E_3^{(+1)}$.
  \item $s_1\, w = w$: the original $V^H$ satisfies $s_1\, v = v$;
    all subsequent projections in blocks~$1$ and~$2$ end with $P_1$,
    restoring $E_1^{(+1)}$; and $P_3$ uses $s_3$ which commutes with
    $s_1$ ($|1 - 3| = 2$), preserving $E_1^{(+1)}$.
\end{enumerate}

Suppose $P_2$ annihilates the $E_4^{(+1)}$ component of some
nonzero~$w$ in this image.  Decompose $w = w_+ + w_-$ in
$E_4^{(\pm 1)}$.  If $(1 + s_2)\, w_+ = 0$, then since
$s_2$ commutes with $s_4$, we obtain $(1 + s_4)(1 + s_2)\, w = 0$.
But $w \in E_1^{(+1)} \cap E_3^{(+1)}$, so
Lemma~\ref{lem:k4} forces $w = 0$---a contradiction.

Therefore $P_2$ is injective on the transported $V^H$ image, and
no rank drop occurs.
\end{proof}


%% ====================================================================
\section{Obstruction for $k \ge 6$}
%% ====================================================================

\begin{remark}[Why the argument does not extend to $k = 6$]
\label{rem:k6-obstruction}
For even block $k = 6$, the analogous contraction would require
showing that
\[
  E_1^{(+1)} \cap E_5^{(+1)} \cap
  \ker\!\left(\frac{(1 + s_6)(1 + s_4)}{4}\right) = \{0\}.
\]
The proof of Lemma~\ref{lem:k4} used two adjacent disjointness
applications: first $(s_1, s_2)$ to kill the $E_4^{(+1)}$ component,
then $(s_3, s_4)$ to kill the remainder.

At $k = 6$, Step~1 would decompose in $E_6^{(\pm 1)}$ eigenspaces.
Since $s_4$ commutes with $s_6$ ($|4 - 6| = 2$), projecting
$(1 + s_6)(1 + s_4)\, v = 0$ to $E_6^{(+1)}$ gives
$(1 + s_4)\, v_+ = 0$, hence $v_+ \in E_4^{(-1)}$.  Combined with
$s_1\, v_+ = v_+$, we get $v_+ \in E_1^{(+1)} \cap E_4^{(-1)}$.
But $s_1$ and $s_4$ are \emph{not} adjacent ($|1 - 4| = 3$), so
adjacent disjointness does not apply.  Indeed,
$E_1^{(+1)} \cap E_4^{(-1)} \ne \{0\}$ in general (e.g., in the
standard representation of $\Sym_7$).

The natural repair would be to use $s_3$ (adjacent to $s_4$) instead
of~$s_1$.  This would require $v \in E_3^{(+1)}$ at the gap step.
However, the cascade through blocks~$3$ and~$4$ includes $P_4$,
which uses $s_4$ adjacent to $s_3$ ($|3 - 4| = 1$), and this breaks
the $E_3^{(+1)}$ eigenspace condition.  The transported $V^H$ image
at the $k = 6$ gap step lies in $E_5^{(+1)} \cap E_1^{(+1)}$ but
\emph{not} in $E_3^{(+1)}$.

Computational verification confirms that the gap is safe through
$\Sym_8$, but a structural proof for $k \ge 6$ remains open.
\end{remark}


%% ====================================================================
\section{Computational verification}
%% ====================================================================

All claims have been verified computationally using exact arithmetic
over~$\mathbb{Q}$ (rational arithmetic in Python/SageMath) via
Young's seminormal form.

\begin{proposition}[Computational verification]\label{prop:comp}
The following have been verified:
\begin{enumerate}
  \item \textbf{Lemma~\ref{lem:k4}:}
    $E_1^{(+1)} \cap E_3^{(+1)} \cap \ker(Q') = \{0\}$,
    where $Q' = (1 + s_4)(1 + s_2)/4$, holds for all irreducible
    representations $V_\lambda$ of\/ $\Sym_n$ with $n = 5, 6, 7$.

  \item \textbf{$V^H$ version:}
    $V_\lambda^H \cap \ker(Q') = \{0\}$ holds for all irreps
    of\/ $\Sym_n$, $n = 5, 6, 7, 8$.

  \item \textbf{Transported $V^H$:} The staircase-transported image
    at the gap step intersects $\ker(Q')$ trivially for all irreps
    through $\Sym_7$.
\end{enumerate}
\end{proposition}

By parabolic reduction, even block $k = 4$ depends only on $\Sym_5$
irreps.  Since the algebraic proof of Lemma~\ref{lem:k4} applies to
\emph{any} $\Sym_n$-representation, the result is unconditional.


%% ====================================================================
\section{Summary of proof status}
%% ====================================================================

\begin{center}
\begin{tabular}{lll}
\toprule
\textbf{Component} & \textbf{Status} & \textbf{Scope} \\
\midrule
Within-block cascade & \textbf{Proved} & All $n$ \\
Block $2$ start & \textbf{Proved} & All $n$ \\
Odd block starts & \textbf{Proved} & All $n$ \\
Even block $k = 4$ & \textbf{Proved} & All $n \ge 5$ \\
\midrule
Even block $k = 6$ & Verified & $n \le 8$ \\
Even block $k = 8$ & Verified & $n \le 16$ \\
Even block $k \ge 6$ general & \textbf{[Gap]} & \\
\bottomrule
\end{tabular}
\end{center}

\medskip
With the $k = 4$ closure, the Rank Isolation Lemma is proved
unconditionally for all $n \le 6$ (blocks $\le 5$).  For $n \le 16$,
the remaining even-block gaps ($k = 6, 8, \ldots, 14$) are closed by
computational verification.  A uniform proof for all even $k$ remains
the principal open problem.

\end{document}
